\section{Introduction}
\begin{itemize}
	\item why ligo needs high power
	\item brief history of LIGO laser 
	\item mention noise problems leading to wanting new amplifier
\end{itemize}
\section{Requirements}
\begin{itemize}
	\item power requirements
	\item long term stability requirement
	\item noise requirements (pointing+frequency+amplitude, only measured frequency and amplitude)
\end{itemize}
\section{OEM solid-state Amplifier}
\begin{itemize}
	\item why are they better than an HPO in principle
	\item how they work
	\item neovan 4S-HP
\end{itemize}
\section{Set-up and Layout}
\begin{itemize}
	\item diagram of layout
	\item description of layout
	\item use 2 amplifiers and a PMC to get high power gaussian beam
	\item flow rates and beam quality of amplifiers
\end{itemize}
\begin{figure}
	\centering
	\includegraphics[width=\linewidth]{figures/E1900168-v3_O4laser_TNT_TableLayout_AOM}
	\caption{e1900168. Layout placeholder}
	\label{fig:e1900168-v3o4lasertnttablelayoutaom}
\end{figure}

\section{Frequency Noise}
Need to put together the measurements for frequency noise.
\begin{figure}
	\centering
	\includegraphics[width=0.7\linewidth]{figures/PMC_open_loop_transfer_function}
	\caption{Open loop gain of frequency stabilisation loop}
	\label{fig:pmcopenlooptransferfunction}
\end{figure}
\section{Long Term Power Results}
In touch with Matt about his, I would like 1 month or something like that (as long as possible)
\begin{figure}
	\centering
	\includegraphics[width=0.7\linewidth]{figures/mode_scan_comparisons}
	\caption{Flow rates affecting the thermal lensing, resulting in increased ellipticity}
	\label{fig:modescancomparisons}
\end{figure}

\section{Conclusion}
summary of how much power for how long, and how it could be improved slightly with mode matching
\chapter{Amplitude Stabilisation for High Power Laser at LIGO}
\section{Introduction + Requirements}
\begin{itemize}
	\item shot noise requirement
	\item radiation pressure noise setting noise requirement
	\item inner+outer loop design
	\item could include calculation for mismatches in the IFO? Point absorbers etc.
\end{itemize}
\section{Free Running Noise and Noise Hunting}
\begin{itemize}
	\item free running noise of the laser
	\item what dominates the laser noise? NPRO
	\item noise sources such as mounts 
	\item noise from PMC offset not being tuned
	\item 4Hz noise from Beckhoff
\end{itemize}
\section{AOM Stabilisation}
\begin{itemize}
	\item description of AOM
	\item AOM TF
	\item AOM loop TF
	\item plot of AOM stabilisation working
\end{itemize}
\begin{figure}
	\centering
	\includegraphics[width=0.7\linewidth]{figures/AOM_stabilisation_noise}
	\caption{Free running and stabilised noise of laser}
	\label{fig:aomstabilisationnoise}
\end{figure}
\begin{figure}
	\centering
	\includegraphics[width=0.7\linewidth]{figures/AOM_stabilisation_TF}
	\caption{Transfer functions for parts of the AOM servo}
	\label{fig:aomstabilisationtf}
\end{figure}
\section{Current Shunt Stabilisation}
\begin{itemize}
	\item description of CS
	\item CS TF
	\item CS loop TF
	\item plot of CS stabilisation working
\end{itemize}
\begin{figure}
	\centering
	\includegraphics[width=0.7\linewidth]{figures/current_shunt_tfs}
	\caption{current shunt input to light power transfer function for different diode currents}
	\label{fig:currentshunttfs}
\end{figure}
\begin{figure}
	\centering
	\includegraphics[width=0.7\linewidth]{figures/dc_power_all_pd}
	\caption{DC behaviour of current shunt}
	\label{fig:dcpowerallpd}
\end{figure}
\section{Comparison of Two stabilisation Schemes}
\begin{itemize}
	\item why I think the AOM is a better choice even though they both work
\end{itemize}
\section{Conclusion}